\chapter{Introduction and Motivation}
\label{ch:intro}

近年來，人工智慧(Artificial Intelligence, AI)與雲端數據中心的快速發展，
使全球對高效能運算(High-Performance Computing, HPC)與高速資料傳輸的需求持續提升。
然而，傳統數據中心與晶片內部所採用的電互連(Electrical Interconnects)架構，
受到金屬導線歐姆損耗所造成的焦耳熱(Joule Heating)、訊號衰減以及頻寬限制等問題影響，
逐漸面臨功耗與資料傳輸能力上的瓶頸。

為了突破上述限制，矽光子(Silicon Photonics)技術利用光訊號進行資料傳輸，
具備高頻寬、低延遲與低功耗等優勢，
因此被視為未來高速晶片互連與光子積體電路(Photonic Integrated Circuits, PICs)的重要技術之一。
隨著矽光子元件逐漸朝向高密度整合與小型化發展，
元件的幾何結構亦變得更加複雜，
使得如何在有限的設計區域內尋找具有良好光學性能的結構，
成為矽光子元件設計中的重要研究課題。

在光子積體電路中，彎曲波導(Bend Waveguide)、交叉波導(Waveguide Crossing)
以及模態分波解多工器(Mode-Division Demultiplexer, DEMUX)
皆為常見且重要的光學元件。
彎曲波導可用於改變光訊號的傳播方向，
交叉波導則可使不同光路在有限晶片面積內相互交會，
而 DEMUX 則可將不同傳播模態分離至對應的輸出通道，
以支援模式分波多工(Mode-Division Multiplexing, MDM)系統。

然而，當元件尺寸逐漸縮小至微米甚至奈米尺度時，
光場與幾何結構之間的交互作用變得更加複雜。
彎曲波導可能因光場偏移與散射而產生額外損耗；
交叉波導則需同時抑制插入損耗(Insertion Loss)與通道間串擾(Crosstalk)；
而 DEMUX 更需要在有限的設計空間內實現不同模態之有效分離。
因此，傳統依賴人工經驗調整幾何參數的方法，
在面對高自由度與複雜光學響應的結構時，
往往難以有效探索整體設計空間。

目前矽光子元件的最佳化方法中，
其中一類為基於光學經驗與解析模型的啟發式試誤法(Heuristic Trial-and-Error Approach)，
此方法雖然適用於結構較為單純的標準光學元件，
但當設計自由度增加時，
其設計效率與可探索的結構範圍會受到限制。

另一類則為目前廣泛應用於奈米光子反向設計中的
基於梯度之伴隨方法(Adjoint Method)。
此方法可透過一次正向電磁模擬(Forward Simulation)
與一次伴隨模擬計算設計變數之梯度，
並依據梯度資訊逐步更新結構，
因此能有效處理具有大量設計變數的拓撲最佳化問題。
然而，由於其本質上仍屬於局部梯度搜尋方法，
最佳化結果會受到初始結構與最佳化路徑影響，
並可能收斂至不同的局部最佳解(Local Optimum)。

為了探索不同於梯度式最佳化的搜尋方式，
近年來已有研究將因數分解機(Factorization Machine, FM)
與退火演算法(Annealing Algorithm)結合，
利用 FM 建立設計變數與目標函數之間的代理模型(Surrogate Model)，
並將其轉換為二元二次最佳化
(Quadratic Unconstrained Binary Optimization, QUBO)問題，
再透過量子退火(Quantum Annealing, QA)或模擬退火(Simulated Annealing, SA)
搜尋具有較佳目標函數值的候選解。
此類方法已被應用於光學薄膜與熱輻射結構等設計問題，
顯示其具有進行全域設計空間搜尋的潛力。

受到上述研究啟發，
本研究將 Factorization Machine 與 Simulated Annealing 結合，
建立 Factorization Machine with Simulated Annealing (FMSA) 反向設計架構，
並進一步將其應用於二維矽光子元件之結構最佳化。

然而，相較於一維薄膜結構，
二維矽光子元件通常具有更高維度的設計空間。
若直接將設計區域離散為 $N\times N$ 個二元像素，
則其可能結構數量可表示為

\begin{equation}
2^{N\times N}.
\end{equation}

隨著解析度提升，
設計變數數量將快速增加，
使搜尋空間呈指數成長。
此外，在 FMSA 架構中，
每一組候選結構皆需透過電磁模擬取得其光學性能，
因此可取得的訓練樣本數量通常遠小於完整設計空間。
當設計維度持續增加時，
有限的訓練樣本將使 Factorization Machine 更難準確描述
設計變數與光學性能之間的關係，
同時也增加退火演算法搜尋 QUBO 問題的複雜度。

因此，對於 FMSA-based photonic inverse design 而言，
\textbf{如何建立有效的結構編碼方式，
以較少的設計變數表示具有足夠幾何自由度的光學結構，
是降低最佳化複雜度的重要問題。}

基於上述研究動機，
本研究的主要目標為
\textbf{尋找適合 FMSA 光子反向設計的高效結構編碼方法}。
透過結構編碼，
將原始高維像素空間轉換為較低維度或較具結構性的表示空間，
使有限的模擬樣本能更有效地描述設計結構與光學性能之間的關係，
並提升 Factorization Machine 與退火搜尋的最佳化效率。

為此，本研究導入三種不同的結構編碼策略，
包括高斯上採樣(Gaussian Upsampling, GU)、
二維傅立葉編碼(2D Fourier Basis Encoding)
以及 $\beta$-變分自編碼器
($\beta$-Variational Autoencoder, $\beta$-VAE)。

Gaussian Upsampling 透過低解析度二元變數經由平滑與投影方式生成高解析度結構，
藉此降低最佳化變數數量；
2D Fourier Basis Encoding 則利用有限數量的頻率基底描述二維幾何形狀，
以頻域係數取代直接像素化表示；
而 $\beta$-VAE 則透過生成式模型學習高維結構資料中的低維潛在空間(Latent Space)，
使 FMSA 能夠直接於潛在變數中進行搜尋。

三種方法皆試圖在降低最佳化維度的同時，
保留足夠的結構表示能力，
但其在幾何自由度、代理模型學習難度與最佳化可擴展性方面具有不同特性。
因此，本研究將進一步比較不同編碼策略於二維矽光子反向設計中的表現，
並探討其適用範圍。

為驗證所提出方法之有效性，
本研究分別針對 90 度彎曲波導、
交叉波導以及模態分波解多工器進行反向設計，
並從光學性能、最佳化行為以及設計變數增加時的可擴展性等面向，
分析不同結構編碼策略與 FMSA 架構之最佳化能力。

本論文之主要貢獻如下：

\begin{itemize}

\item \textbf{建立 FMSA 二維矽光子反向設計架構：}
本研究將 Factorization Machine 與 Simulated Annealing 結合，
並應用於二維矽光子元件之結構最佳化，
利用代理模型與退火搜尋進行非梯度式設計空間探索。

\item \textbf{導入並比較三種結構編碼策略：}
本研究分別採用 Gaussian Upsampling、
2D Fourier Basis Encoding 與 $\beta$-VAE
作為二維矽光子結構表示方法，
探討不同編碼方式在降低設計維度、
保留幾何表示能力以及代理模型學習效率方面之差異。

\item \textbf{分析不同編碼策略與設計維度對 FMSA 最佳化表現之影響：}
本研究將所提出之方法應用於 90 度彎曲波導、
交叉波導以及模態分波解多工器，
並透過不同設計變數數量之最佳化結果，
分析設計空間維度對 Factorization Machine 學習能力、
退火搜尋行為及最終光學性能之影響，
進一步探討不同編碼策略於二維矽光子反向設計中的適用性。

\end{itemize}